\documentclass[12pt]{article}

\usepackage[english]{babel}
\usepackage[utf8x]{inputenc}
\usepackage{amsmath}
\usepackage{graphicx}
\usepackage[colorinlistoftodos]{todonotes}
\usepackage{listings}
\usepackage{glossaries}
\usepackage{placeins}
\usepackage{fixltx2e}
\usepackage{pdfpages}
\usepackage{lastpage}
\usepackage{enumitem}
\usepackage{xcolor}
\usepackage{scrpage2}
\usepackage{scrtime}
\usepackage{parskip}
\usepackage{hyperref} 

\clearscrheadfoot
\pagestyle{scrheadings}
\usepackage[
top    = 2.5cm,
bottom = 3cm,
left   = 3cm,
right  = 3cm]{geometry}
\setcounter{secnumdepth}{4}
\title{Diploma Thesis}


\author{RoboNav}
\date{\today}


%Includes
%!TEX root=../document.tex

\usepackage[in]{fullpage}
% Fontencoding for possible copy&paste out of PDF
\usepackage[T1]{fontenc}
\usepackage[utf8]{inputenc}
\usepackage[ngerman]{babel}
\usepackage{graphicx} 
\usepackage{textcomp}
\usepackage{sectsty}
\usepackage{caption}
\usepackage{listings}
\usepackage{array}
\usepackage{colortbl}
\usepackage{footmisc}
\usepackage{fancyhdr}
\usepackage{ccicons}
\usepackage{suffix}
\usepackage{multirow}
\usepackage{tabularx}
\usepackage{listings}
\usepackage{color}
\usepackage{url}
\usepackage[dvipsnames]{xcolor}
\usepackage[longnamesfirst,nonamebreak]{natbib}
\usepackage[headsep=1cm,headheight=3cm,hmargin=2cm,vmargin=2.5cm]{geometry}
\usepackage[nolist]{acronym}

% Definitions for Textcolor
\usepackage{color}
\definecolor{listings}{rgb}{0.96, 0.96, 0.96}
\definecolor{update}{rgb}{1, 0.8, 0.8}
\definecolor{config}{rgb}{0.8, 1, 0.8}
\definecolor{gray}{rgb}{0.4,0.4,0.4}
\definecolor{darkblue}{rgb}{0.0,0.0,0.6}
\definecolor{cyan}{rgb}{0.0,0.6,0.6}

% Java Syntaxhighligthning
% strings
\definecolor{javared}{rgb}{0.6,0,0}
% comments
\definecolor{javagreen}{rgb}{0.25,0.5,0.35}
% keywords
\definecolor{javapurple}{rgb}{0.5,0,0.35}
% javadoc
\definecolor{javadocblue}{rgb}{0.25,0.35,0.75}

\lstset{
	basicstyle=\ttfamily\small,
	keywordstyle=\bfseries\color[rgb]{0.496,0.000,0.332},
	commentstyle=\color[rgb]{0.246,0.496,0.371},
	stringstyle=\color[rgb]{0.164,0.000,0.996},
	tabsize=4,
	breaklines=true,
	numbers=left,
	numberstyle=\tiny\color{black},
	stepnumber=5,
	numbersep=10pt,
	numberstyle=\tiny,
	captionpos=b,
	xleftmargin=1cm,
	showspaces=false,
	showstringspaces=false,
	basewidth={0.53em,0.45em},
	frame=single,
	xleftmargin=1cm,
	basicstyle=\scriptsize,
}


\lstdefinestyle{Java}{
	language=Java,
	keywordstyle=\color{javapurple}\bfseries,
	stringstyle=\color{javared},
	commentstyle=\color{javagreen},
	morecomment=[s][\color{javadocblue}]{/**}{*/},
}

\lstdefinestyle{XML}{
	language=XML,
	basicstyle=\ttfamily,
	columns=fullflexible,
	commentstyle=\color{gray}\upshape
}

\lstdefinelanguage{XML}
{
	morestring=[b]",
	morestring=[s]{>}{<},
	morecomment=[s]{<?}{?>},
	stringstyle=\color{black},
	identifierstyle=\color{darkblue},
	keywordstyle=\color{cyan},
	% list your attributes here
	morekeywords={xmlns,version,type}
}

\usepackage[
	colorlinks,
	citecolor=black,
	filecolor=black,
	linkcolor=black,
	urlcolor=black,
	linktoc=all
]{hyperref}

\makeglossaries
\newglossaryentry{json} {name=JSON, description={Java Script Object Notation}}
\newglossaryentry{GPS} {name=GPS, description={Global Positioning System }}
\newglossaryentry{NDGPS} {name=NDGPS, description={National Differential Global Positioning System }}
\newglossaryentry{WADGPS} {name=WADGPS, description={Wide Area Differential Global Positioning System }}
\newglossaryentry{QR} {name=QR, description={Quick Response (Code) }}
\newglossaryentry{RNCP} {name=RNCP, description={RoboNav Communication Protocol}}
\newglossaryentry{TGM} {name=TGM, description={Technologisches Gewerbe Museum / Institute of technology}}
\newglossaryentry{API} {name=API, description={Application Programming Interface}}
\newglossaryentry{JNI} {name=JNI, description={Java Native Interface}}
\newglossaryentry{WiFi} {name=WiFi, description={Wireless Local Area Network}}
\newglossaryentry{DHCP} {name=DHCP, description={Dynamic Host Configuration Protocol}}
\newglossaryentry{IP} {name=IP, description={Internet Protocol}}
\newglossaryentry{MOV} {name=MOV, description={QuickTime Movie (file extension)}}
\newglossaryentry{MPEG4} {name=MPEG4, description={Moving Picture Experts Group 4 (file extension)}}
\newglossaryentry{MJPEG} {name=MJPEG, description={Motion Joint Photographic Experts Group (file extension)}}
\newglossaryentry{JPEG} {name=JPEG, description={Joint Photographic Experts Group (file extension)}}
\newglossaryentry{CPU} {name=CPU, description={Central Processing Unit}}
\newglossaryentry{GHz} {name=GHz, description={Gigahertz (thousands of MHz)}}
\newglossaryentry{ROS} {name=ROS, description={Robot Operationg System}}
\newglossaryentry{GNUGPL} {name=GNU GPL, description={Gnu's Not Unix General Public License}}
\newglossaryentry{BSD} {name=BSD, description={Berkeley Software Distribution}}
\newglossaryentry{GUI} {name=GUI, description={Graphical User Interface}}
\newglossaryentry{DFT} {name=DFT, description={Discrete Fourier Transform}}
\newglossaryentry{RGB} {name=RGB, description={Red Green Blue}}
\newglossaryentry{UML} {name=UML, description={Unified Modeling Language}}
\newglossaryentry{MRDS} {name=MRDS, description={Microsoft Robotics Developer Studio}}
\newglossaryentry{IDE} {name=IDE, description={Integrated Development Environment}}
\newglossaryentry{FXML} {name=FXML, description={JavaFX - Extensible Markup Language file}}
\newglossaryentry{XML} {name=XML, description={Extensible Markup Language}}
\newglossaryentry{MVC} {name=MVC, description={Model View Control}}
\newglossaryentry{CSS} {name=CSS, description={Cascading Style Sheet}}

\renewcommand*\glspostdescription{\dotfill}




\begin{document}


\begin{titlepage}
\begin{center}

% Additional to first page
\includegraphics[width=1.0\textwidth]{images/robonavlogo}\\

\LARGE TGM - HTBLuVA Wien XX \\ IT Department  \\[1.5cm]

% Title
\rule{1.0\textwidth}{1mm}
{ \huge \bfseries \\[0.4cm]  \huge Diploma Thesis \\ \LARGE RoboNav \\[0.4cm] }

\rule{1.0\textwidth}{1mm}



\noindent 


\vfill

% Bottom of the page
{\small Version: \today ~at  \thistime    }
\end{center}

%\end{center}
\end{titlepage}


%HEADER AND FOOTER
\pagenumbering{Roman}
\ohead{\headmark}
\ifoot{© RoboNav - 2014/15}
\ofoot{\pagemark}

\newpage %----------------------------------------------------------------------------------------------
{\small\color{white}.}
\vspace{-0.7cm}
\tableofcontents

\newpage %----------------------------------------------------------------------------------------------
% Affirmation
\includepdf[scale=1.2,pages=-]{documents/soor}
\newpage
\section*{Abstract}
\cfoot{Hannah Siegel}

Nowadays robotics has become one of the most crucial foundations of a day-to-day business in nearly every industrial environment. As a result of the cost-efficiency and the overall benefits of autonomous systems, robots are in high demand in this area and their technologies are put under constant improvement. It is the strong presence of robots which distinguishes factories in industrialized countries from their competitors in regions with low-cost production possibilities. Nevertheless in the field of mobile robotics, used mostly for logistic purposes, we have come to witness an extreme unused potential. Mobile systems are not always able to detect and therefore react to new changes or obstacles in their surrounding areas and this increases the risk of damaging material or – worse – of hurting humans.  Mostly these security concerns may lead to the fact that the path-planning possibilities of mobile robots are limited to only pre-defined and hardcoded pathways which excessively scaled down the flexibility of these systems.
The project - \textit{RoboNav} - aims to evaluate if a system with better positioning and more flexible navigation can be established if external image processing units deliver additional real-time information from a bird's eye perspective. 
In order to simulate these technical procedures in a school project, the overall idea had to be scaled down to simpler hardware and software environments while constantly trying to stick to the theoretical industrial environment and its challenges. The proof of concept comprises four entities: a mobile execution device, the \textit{Robotino® v3} from Festo AG \& co. KO, the external image capturing unit (\textit{RoboNav Sight}),  represented by an overhead Android smart phone with a streaming app, an intuitive graphical user interface (\textit{RoboNav Overview}) and in the center of the system the control application (\textit{RoboNav Control}), which contains the important aspects of \textit{RoboNav}. The responsibilities of \textit{RoboNav Control} are numerous: It acts as a communication middleware for all the other components, handles the image processing and information gathering structures and has to provide path-planning and path-driving functionalities. Additionally, a asynchronous communication link has been developed, \textit{RoboNav Communication Protocol} (\textit{RNCP}). The important part of improving the navigation is done by position comparison between the internal position of the robot's odometry and the correct position as seen from an overhead camera. The project is a proof of concept, but the outcome of this work can be easily transformed onto other systems with practically no changes to be done.

\newpage
\section*{Kurzfassung}
\cfoot{Vennesa Belinić}

Robotik ist heute eine der wichtigsten Technologien im industriellen Umfeld. Durch die Kosteneffizienz und den Gesamtvorteil autonomer Robotersysteme sind diese und ihre Technologien sehr gefordert und stehen unter ständiger Entwicklung. Es ist der starke Einsatz und die Präsenz von Robotern, welche Unternehmen in Industrieländern von anderen mit kostengünstigeren Produktionsmöglichkeiten unterscheidet.\\
Mobile Robotersysteme werden meist für logistische Zwecke eingesetzt.
Dieses sind nicht immer in der Lage auf alle Änderungen ihrer Umgebung dynamisch zu reagieren. Dadurch können nicht nur Schäden an anderen Geräten zustande kommen, sondern auch Menschen die sich in deren Umgebung aufhalten in Gefahr gebracht werden.
Aufgrund dieser Einschränkung und auch wegen Sicherheitsbestimmungen können sich die meisten mobilen Robotersysteme nur auf vordefinierten Bahnen oder Arealen bewegen. Dies schränkt die Flexibilität des gesamten Systems ein.
Das Projekt \textit{RoboNav} verfolgt das Ziel, zu evaluieren ob eine flexiblere und genauere Navigation mit externen Bildverarbeitungssystemen aus Vogelperspektive erreicht werden kann.
\textit{RoboNav} hat die Systemumgebung so weit vereinfacht um dieses Szenario in diesem Schulprojekt nachzubauen. Trotzdem wird versucht noch den theoretischen industriellen Anforderungen zu entsprechen.
Das Proof-of-Concept umfasst vier Einheiten: das mobile Robotersystem \textit{ Robotino® v3} von Festo, die externe Bildaufnahme (\textit{RoboNav Sight}) mittels eins Android Smartphone mit einer Streaming-Applikation, eine Interaktionsmöglichkeit mit dem System durch eine intuitive graphische Oberflächen (\textit{RoboNav Overview}) und die Kontrollsoftware des Systems (\textit{RoboNav Control})  welche die wichtigsten Aspekte dieses Projektes beinhaltet. \textit{RoboNav Control} verhält sich wie eine Middleware und verbindet alle anderen Komponenten miteinander, teils durch ein selbst definiertes Kommunikationsprotokoll (\textit{RNCP}). Die Kontrollsoftware beinhaltet auch die Bildbeschaffung und -verabeitung und stellt Pfadplanungs- und Pfadbefahrungsfunktionalitäten zu Verfügung. Durch den Abgleich der tatsächlichen Position mit der internen vom Roboter berechneten ist der wichtigste Part des Projekts abgedeckt.
Das Projekt ist als Machbarkeitsstudie zu betrachten, trotzdem sollte das Resultat eine Software sein die sich ohne große Designänderungen in komplexere, andere Systeme einbauen lässt.

\newpage
\section*{Acknowledgements}
\cfoot{}

We want to thank Festo AG \& Co. KG for sponsoring the \textit{Robotino v3}, PRIA (Practical Robots Institute Austria) for contributing their expertise and our tutors Erhard List BSc. and Michael Borko BSc. for their priceless advice, their constant encouragement and their support in terms of project management related know-how.


\label{pageRomanEnd}

\newpage %----------------------------------------------------------------------------------------------
\pagenumbering{arabic}
\ofoot{\pagemark}

\section{Introduction}
% Basic Introduction, Goal of the Project and very short description of the system (see section ...)
\label{sec:introduction}

\subsection{Navigation with mobile robots}
\label{subsec:stateoftheart}
%\input{}

\subsection{The Project "RoboNav"}
\label{subsec:project}
%\input{}

\subsection{Project Management}
\label{subsec:pm}
%\input{}

\newpage %----------------------------------------------------------------------------------------------

\section{System design}
\label{sec:systemdesign}
%\input{}

\subsection{Domain model}
\label{subsec:domainmodel}
%\input{}

\newpage
\subsection{Component Overview}
\label{subsec:components}
%\input{}

\subsubsection{Robotino v3}
\label{subsubsec:robotinov3}
%\input{}

\subsubsection{Overhead Camera (an Android Smartphone)}
\label{subsubsec:ipwebcam}
%\input{}

\newpage %----------------------------------------------------------------------------------------------
\section{RoboNav Communication Protocol (RNCP)}
\label{sec:rncp}
%\input{}
\newpage

\subsection{Protocol Specification}
\label{subsec:protocolspec}
%\input{}

\newpage
\subsection{Implementation of the Transmission}
\label{subsec:impltransmission}
%\input{}

\subsection{Using RNCP}
\label{subsec:unsingrncp}
%\input{}

\newpage %----------------------------------------------------------------------------------------------
\section{RoboNav Control}
\label{sec:controlsoftware}
%\input{}

\subsection{Controlling the interaction between the components}
\label{subsec:interaction}
%\input{}

\newpage
\subsection{Path planning}
\label{subsec:pathplanning}
%\input{}

\newpage %----------------------------------------------------------------------------------------------
\section{Image Retrieval and Processing}
\label{sec:imageproc}
%\input{}

\subsection{Retrieving an Image}
\label{subsec:imgretrieving}
%\input{}

\subsection{Analysing Images}
\label{subsec:imganalysing}
%\input{}

\newpage
\subsection{Object Classification}
\label{subsec:imgclassification}
%\input{}

\newpage %----------------------------------------------------------------------------------------------
\section{Interacting with the Robotino v3}
\label{sec:robotinov3}
%\input{}

%JNI
\subsection{Java API}
\label{subsec:robotinojavaapi}
%\input{}

\subsection{Communication with the Robotino}
\label{subsec:communicationRobotino}
%\input{}

 \subsection{Path-driving functionality}
\label{subsec:pathrivingr3}
\cfoot{Hannah Siegel}
One important aspect of the \textit{RoboNav} project is the path-planing based on information gathered from the overhead camera (see also section \ref{subsec:pathplanning}). The robot can make a request to the \textit{RoboNav Control} with coordinates it want to drive to or a request can be send over the \textit{RoboNav Overview}. This section discusses the path-driving functionality in detail.
%The robot should be able to move to a destination a corresponding path-driving function is needed as well. \\
%Once the target of the robot is known and the path planning algorithm (section \ref{subsec:pathplanning}) has calculated the fastest path avoiding obstacles the path will be driven. \\
%The Path-driver's duty is it to let a mobile vehicle - in the case of this project this is exclusively the \textit{Robotino v3} - drive this path very accurately.
\subsubsection{High-level path-driving of the Robonav Control with Java}
\label{subsubsec:pathdrivinghighlevel}
Whenever the robot has to move to a certain point, a \texttt{DrivablePath} object, which is implementing the \texttt{Path} interface, is generated by the path planning algorithm. A  \texttt{DrivablePath}, generated by the path-planning software, acts as the component containing the right information needed for the path driving. \\
In order to enable the \textit{Robotino v3} as an actual device to move on the course of the \texttt{DrivablePath}, a class called \texttt{RobotinoPathDriver} was implemented which extends the super class \texttt{PathDriver}.\insertpicture{images/uml_pathdriver}{Path-driving \gls{UML}}{(selfmade)}{fig:pathdriveruml}{1.0}
The program flow is as follows:
\begin{enumerate}
\item A path-request was either done by the robot (see section \ref{subsec:robotinoview} for details) or over the \textit{Robonav Overview}.
\item When requested, a new path is generated by the path planning and stored in an \texttt{DrivablePath} object. This happens in the the \textit{RoboNav Control}. 
\item The \textit{RoboNav Control} charges a \texttt{PathDriver} with the driving logic, by handling the \texttt{DrivablePath} to the \texttt{PathDriver}. The \texttt{RobotinoPathDriver} class encapsulates the path-driving special for the \textit{Robotino v3}. For other robot models, this class must be implemented accordingly. The usage of the Path driver can can be seen in Listing  \ref{ExternalPathDriving}.
\item The \textit{Robotino} brings a powerful Path-following software with its \textit{RobotinoView} Software. \textit{RoboNav} makes use of this by transferring the path information from the \texttt{RobotinoPathDriver} to the \textit{Robotino} using text-files which are read by the \textit{RobotinoView} path-following software. \\
This special software is closely discussed in section \ref{subsubsec:pathdriverRobotino} and the transferring of the files is discussed section \ref{subsec:communicationRobotino}.
%The \texttt{PathDriver} will let the hardware drive a pathway. In the case of \textit{RoboNav} and the \textit{Robotino v3} this has been done using \textit{RobotinoView} since this was easier than implementing it in Java. The function is discussed in section \ref{subsubsec:pathdriverRobotino}. For this purpose the \texttt{DrivablePath} must be transformed into a list of strings and the right commands must be sent to the Robotino using the, in section \ref{subsec:communicationRobotino}, explained file transfer methods for writing and reading remote files. The driving logic itself will then be done by the RobotinoView program so there is no need for complicated calculations within the API and because only command-structures are send the network is not getting used too much and therefore the path-driving is not vulnerable from long latency or network errors.
\end{enumerate}
%A few important steps on the side of the \textit{RoboNav Control} are mainly the right casting of a path to a \texttt{String} list and the correct usage of the commands so that the \textit{RobotinoView} program can make use of the information transmitted.\\ 
The \texttt{DrivablePath} contains a list of \texttt{CoordinateVectors} where each item contains the next node (= point the robot should drive to) as explained in section \ref{subsubsec:graphbasedrepresentation}. \\
These points are treated as sub-paths by the path driving algorithm. The mobile robot drives in a straight movement from the current position to the next point in the list without rotating. At the end of each sub-path the robot may then turn to a specific angle. This is especially convenient if the robot has a grabber or some similar part which may obstruct in any way if not driving with the right angle. \\ In Listing \ref{transformationPathDriving} the transformation of a high-level path to a more low level string list can be seen. The method returns a \texttt{String}-list containing quadruples of data in the form described in Listing   \ref{Resultofthetransformation} (containing an counter, the x and y values and the rotation angle called phi of the next sub-path point). Before storing the x and y values into the list, they are multiplied by 1000, since the \textit{Robotino} works with millimetres, but \textit{RoboNav} uses metres. Additionally, the whole path string must have an terminating statement which is defined by for times zero. This termination statement tells the path driver to wait until a new path command has been sent. 
\begin{lstlisting}[caption = Simplified transformation logic, label = transformationPathDriving]
// transforming the points using a lambda expression
this.getAnchorPoints().forEach(p -> pointsAsStringList.add(ai.getAndIncrement()+" "+p.getX()*1000+" "+p.getY()*1000+" "+p.getPhi()) );
    
// adding the terminatory statement
pointsAsStringList.add("0 0 0 0");
\end{lstlisting}
\begin{lstlisting}[caption = Result of the transformation, label = Resultofthetransformation]
1 1000.0 1000.0 0.0
2 2000.0 1000.0 90.0
3 2500.0 -500.0 180.0
0 0 0 0
\end{lstlisting}
\subsubsection{Controlling the path driver's actions} 
After the path is transformed, the task of transmitting the data to the \textit{Robotino} begins. \\ 
The more low level path driving logic uses five files for communication, as defined in Listing  \ref{filesPathDriving}. These files must get used in the right order and with the right values, in order to ensure that the path-driving software on the \textit{Robotino} which is discussed in section executes the right actions. In Figure \ref{fig:statemachine}, the possible states of the path-driving software can be seen: % Files get used because it is the easiest way of communication and the \textit{RobotinoView} application supports reading and writing files in a very good way.  
\\ 
The path-driver always starts in the \texttt{waiting} state. As soon as a new path is available, it parses each sub step of this path and drives it, the robot is then in the \texttt{driving} state. In case of an emergency stop, the robot stops immediately and goes over to the \texttt{blocked} state. Once the emergency stop is finished, the robot drives on. Finally, if a slow stop was triggered, the robot will just stop when parsing the next subpath, instead of being blocked right away. Once a path has been completely parsed and there is no next sub path available, the robot is done moving and is back in the \texttt{waiting} state again. 
\insertpicture{images/statemachine}{Overview of the path-drivers states}{(self-made)}{fig:statemachine}{1.0}
The actions described in Figure \ref{fig:statemachine} are triggered using the following files: \\ 
The \texttt{RoboNavPath.txt} file contains the subpath list like in Listing \ref{Resultofthetransformation}.\\
The \texttt{RoboNavStopEmergency.txt} file is used for an emergency stop. This means, that if the content of the file is set to \texttt{1}, the robot should stop immediately - regardless of his position. \\
The \texttt{RoboNavStopSlowly.txt} file on the other hand sends the signal for a controlled stop at the next sub-path position of the path. For example, this function may be used if the robot should stop driving the current path, so that he can drive another one immediately but first should stop on the end of a sub-path and not just anywhere.\\
Since these two stop mechanisms can be used at the same time, \textit{RoboNav} used these two files instead of one global stop file whereas \texttt{0} would be no stop, \texttt{1} slow stop and \texttt{2} emergency stop. Further, the files are used at two different positions within the path-driver and it is easier to use two files containing only true or false since they can be directly connected into an \texttt{AND-flag} without an if-statement. Finally, the usage of two files comes with the advantage of no possible race-conditions, since \textit{RoboNav} is the only entity modifying these files. 
\\
The \texttt{RoboNavReadyToBeDriven.txt} file acts as an controlling instance for the \textit{RoboNav Control}: if the contents of the file are set to \texttt{1}, the \textit{Robotino} is currently not engaged in any actions and can be used for navigational purposes. Therefore if the robot is presently driving a path, it will be marked as unavailable. Also, it may happen that the transmission of the file or the path-driving function of the \textit{RobotinoView} did not work immediately or that there is no path-driver program running. Therefore, the \texttt{RoboNavReadyToBeDriven.txt} comes with a second benefit: as long as there is no \texttt{0} in the file, the \textit{Robonav Control} knows, that the driving of the path was not successful and transmits the path again. If after five times re-sending the path the path-driver did not react, the \texttt{CanNotDrivePath} 
exception is thrown, which means that there is no \textit{RobotinoView} file running at all. \\ 
Finally, the \texttt{RoboNavNewPathAvailiabe.txt} file is used to trigger the \textit{RobotinoView} program when a new path is available and that this path should be followed immediately.

\newpage
\begin{lstlisting}[caption = Communication Files and their paths on the Robotino, label = filesPathDriving]
1. REMOTE_FILE_PATH = "/home/robotino/RoboNavPath.txt"
2. REMOTE_FILE_STOP_EMERGENCY = "/home/robotino/RoboNavStopEmergency.txt"
3. REMOTE_FILE_STOP_SLOWLY = "/home/robotino/RoboNavStopSlowly.txt"
4. REMOTE_FILE_READYTOBEDRIVEN = "/home/robotino/RoboNavReadyToBeDriven.txt"
5. REMOTE_FILE_NEWPATH = "/home/robotino/RoboNavNewPathAvailiabe.txt"
\end{lstlisting} 
Even through these file structures may seem unnecessary complicated, they enable an ideal information exchange between the \textit{RoboNav Control} path-driver. The file commands are handled by the \texttt{RobotinoPathDriver} class, which internally is using the \texttt{SSHManager} (Listing \ref{InternalPathDriving}) but the \textit{RoboNav Control} only has to perform the actions described in Listing \ref{ExternalPathDriving} if a path has to be transmitted. \\
\textit{RoboNav} is not using \textit{RNCP} for this communication because the file exchange is easily implemented in Java using the Jsch library and on the \textit{Robotino} side using the file processing of \textit{RobotinoView}. More important, the usage of files opens the possibility to leave the \textit{Robotino} system unaltered. There is no additional need for implementing \textit{RNCP} into all the possible programming interfaces of the \textit{Robotino} since every language can work with files.
\begin{lstlisting}[caption = Internal handling of the path driving using the SSHManager (section \ref{subsec:communicationRobotino}), label = InternalPathDriving]
while(sshman.checkFileContent(Robotino3.REMOTE_FILE_READYTOBEDRIVEN,"1")){ 
    // wait until the robotino is ready for navigation
}

// reset new path flag
sshman.writeFile(Robotino3.REMOTE_FILE_NEWPATH,"0");

// delete old path
sshman.clearFile(Robotino3.REMOTE_FILE_PATH);

// write new path (lines = transformed path to string list)
sshman.writeFile(Robotino3.REMOTE_FILE_PATH,lines);

// signal that a new path is availiable
sshman.writeFile(Robotino3.REMOTE_FILE_NEWPATH,"1");
\end{lstlisting}
\begin{lstlisting}[caption = Actions that need to be performed to trigger the path-driving functionality, label = ExternalPathDriving]
RobotinoPathDriver robotinoPathDriver = new RobotinoPathDriver();
robotinoPathDriver.letDriveThePath(driveablePath,robotino);
\end{lstlisting}


\newpage
\subsubsection{Low-level path-driving on the Robotino using RobotinoView}
\label{subsubsec:pathdriverRobotino}
\textit{RoboNav} provides \textit{RobotinoView} code as an example implementation for using the path-driving functionality in own programs. Furthermore, this code is also used from the \textit{RoboNav} project to let the \textit{Robotino} drive the calculated path. 
%For realizing a path driving logic the following program is written using the graphical programming interface for the \textit{Robotino}, the \textit{RobotinoView® 3}.\\
The program contains two steps: the first resets the odometry of the \textit{Robotino} to an initial position. The second step actually does the path-driving. This more complex step is now further discussed and can be seen in Figure \ref{fig:overviewpathdriverrobotino}.
% As it can be seen in image \ref{fig:pathdriverstructure} the structure only contains two steps from which the first is only  resetting the odometry. The second step can be seen in figure \ref{fig:overviewpathdriverrobotino}.
 %\insertpicture{images/structure}{Structure of the main program}{(self-made)}{fig:pathdriverstructure}{0.5}
 \insertpicture{images/newpd}{Overview of the path-driving step}{(self-made)}{fig:overviewpathdriverrobotino}{1.0}
Before starting the program, the \texttt{sendResetPathDriverFiles} method of the \texttt{Robotino3} class should be executed. It resets the communication files to their initial values. Additionally, a path gets set which lets the robot drive ten centimetres to the left and back again so the user can see if all systems work fine. %The program step has small text boxes in the picture. A small overview of the program can be gathered therefore but it will be closer described now: 
\subsubsection{New path logic}
After starting the program (preferably over the \textit{Robotino}'s webinterface), it waits for the new-path-command to be sent. In order to send a new path the \\ \texttt{RoboNavNewPathAvailable.txt} must be set to \texttt{0}. Then the path must be written into the file and then the \texttt{RoboNavNewPathAvailable.txt} must be set to \texttt{1}. After that the code which can be seen in Figure \ref{fig:newpath} takes care of the correct interpretation and handling of the new path transmitted. The code appears a little simplified, for example the \texttt{boolean} values must be multiplied by \texttt{1} in order to cast them to \texttt{float} values. This can still be seen in the overall overview program (Figure \ref{fig:overviewpathdriverrobotino}). 
\insertpicture{images/newpath}{New path and next step logic}{(self-made)}{fig:newpath}{1.0}
The new-path-file gets read continuously (though the file reading block) and the first value gets chosen (the float vector block where index 1 is taken). This value is put into an \texttt{AND-flag} which turns \texttt{true} only once even if the input values stay at \texttt{true}. The \texttt{AND-flag} is needed, because this value triggers the Lua-script afterwards. If it would stay \texttt{true} for a longer timespan it would trigger the new path available information over and over again, but the new path is only available once.  \\
Next the Lua-script takes in some information needed for deciding if and how often the path file should be read.  Whenever the output of this script is \texttt{true}, the next sub-path will be read and put into the path driver. If a new path is announced but the old path is not terminated yet, the script must set it's output to true as long as the incremental counter of the path file is \texttt{1} which means that the path is at the beginning. \\
Listing \ref{leftluascript} shows how the five inputs of the Lua-script can be addressed using the variables \texttt{in1} to \texttt{in5}. The output is set using the variable \texttt{out1}. A global variable \texttt{flankcalled} is set if the new path flag was thrown but the other conditions are not \texttt{true} yet. \\
The inputs are:
\begin{enumerate}[topsep=0pt,itemsep=0ex,partopsep=10ex,parsep=1ex]
  \item \texttt{in1 (terminator)} - The value coming from the terminator script. If this value is \texttt{0 (false)}, the robotino is currently not driving any path.
  \item \texttt{in2 (newpathflag)} - The \texttt{AND flag}'s value. If the new-path-file's value has changed it is set to \texttt{ 1 (true)} but only for one cycle.
  \item \texttt{in3 (robotinoStopped)} - The position-reached-value of the path-driver-block. If the \textit{Robotino} has finished the current sub-path this value will be \texttt{1 (true)} just until a new sub-path is passed to the path driver.
  \item \texttt{in4 (pathStep)} - The current step of the path, which is the incremental number in the path file.
  \item \texttt{in5 (nextSubPath)} - If this value is \texttt{1 (true)} the next sub-path should be driven immediately.
\end{enumerate}
%
%terminator = in1;%
%newpathflag = in2;
%robotinoStopped = in3;
%pathStep = in4;
%nextSubPath = in5;
%
\newpage
\begin{lstlisting}[caption = Lua Script, label = leftluascript]
out1 = nextSubPath;
if ( terminator == 0 ) then
	robotinoStopped = 1;
end
if ( newpathflag == 1 ) then
	flankcalled = 1;
end
if (flankcalled == 1 and robotinoStopped == 1 ) then
	out1 = 1;
	if ( pathStep == 1 ) then 
		out1 = nextSubPath;
		flankcalled = 0;
	end
end
\end{lstlisting}
Basically, the script above is setting the output value to \texttt{true} if:
\begin{itemize}[noitemsep]
\item A new path is available. \\
This happens when the path flag was called at least once and the \textit{Robotino} is not driving any more or has performed a short stop between the sub-paths.
The output is \texttt{true} as long as the path-step is not \texttt{1}, since then the path will be driven normally.
\item The next sub-path step should be driven because the path-driver has finished the previous one.
\end{itemize}
\subsubsection{Path driving}
\insertpicture{images/pathdriver}{Path driver logic (simplified)}{(self-made)}{fig:pathdriver}{0.9}
Once the path-driver gets a new position out of the path-file, the \textit{Robotino} drives to the gathered position and turns into the final pose.
The path-driver block is therefore set to holonomic driving and then holonomic turning. \\
% The path driver always needs the odometry. This is why these values were also passed to it.
Once the path driver has received the new values of a position it is waiting for it's last input (marked with "Neustart" in Figure \ref{fig:pathdriver}) to turn \texttt{true} so it will actually start to drive the sub-path.  \\
As long as the robot is still driving the path, the terminator-Lua-Script has \texttt{1} as output. The terminator-Lua-script checks if the incremental counter equals the stopping condition, which is a \texttt{0}, and if so the robot stops driving. \\
The decision for the next sub-path is done by the \texttt{AND-flag} at the bottom. This value is \texttt{true} if the terminator has not yet terminated the path and the path-driver is currently not driving the previous sub-path any more.\\ Finally, the path driver does all the calculations for the motor velocities for the x and y direction and the rotation and passes them on to the engine. The path driver finally signals if the sub-path endpoint has been reached whenever it's last output (marked with "Pose erreicht" in Figure \ref{fig:pathdriver}) is \texttt{true}.
\subsubsection{Emergency and controlled Stops}
\label{subsubsec:stoppingtherobotino}
In industrial environments it is common, if not even necessary due to regulations, to implement emergency stops mechanisms. \textit{RoboNav} implemets a "soft" version of these very important functions.
\insertpicture{images/stops}{Emergency and slow stop}{(self-made)}{fig:stop}{1.0}
%First of all, the robot is stopping by himself if the sensors report any obstacle which is too near. 
A possibility for sending an emergency stop from \textit{RoboNav Control} has been implemented. 
\textit{RoboNav Control} is able to send such an emergency stop command. The signal is sent via the emergency-stop-file and the stop-logic is situated immediately between the path-driver and the direct control of the engine, so as close to the engine as possible. \\
If the content of the emergency-file is \texttt{1}, the negated value will be \texttt{0}. When the velocity values which the path-driver passes on to the driving logic of the \textit{Robotino} are multiplied with \texttt{0} they are \texttt{0}, which means that the robot won't drive. Of course, it works the same way with the terminating value as well. \\
Normally, an \textit{real} emergency stop implementation would need an restart of the system, an new request of the path or setting an dedicated start command but this type of simulated stop was enough for \textit{RoboNav}'s overall purposes. \\\\
The slow-stop on the other hand causes the robot to finish the current sub-path and then stop. The logic for stopping at the next subpath causes the decision of the next subpath to delay until the slow-stop-signal is deactivated.








\subsection{Integration of the System into Robotino View}
\label{subsec:robotinoview}
%\input{}

\newpage %----------------------------------------------------------------------------------------------
\subsection{Adjustment of the coordinate systems}
\label{subsec:coordinatesystems}
%\input{}


\newpage %----------------------------------------------------------------------------------------------
\section{RoboNav Overview}
\label{sec:gui}
\cfoot{Dominik Scholz}

\textit{RoboNav Overview} is used for visualizing the \textit{RoboNav} system. Although \textit{RoboNav} is modular as described in section \ref{subsec:domainmodel} and therefore the use of \textit{RoboNav Overview} is not mandatory, it is a basic implementation of the \textit{RoboNav \gls{API}} thus giving a visually appealing example of a range of functions the \gls{API} offers. \\

\subsection{Design of RoboNav Overview}

The graphical user interface framework chosen for \textit{RoboNav Overview} was JavaFX, the state-of-the-art framework for rich client applications in Java. Since Java 8 it is included in the Java SE Runtime Environment making it easy to reach a broad user base. JavaFX is designed to encourage the Model-View-Control pattern by providing a \gls{XML}-based way to describe the user interface additionally to the approach of building it with program code. These \gls{XML}-files, called \gls{FXML}, are style-able with a derived version of the \gls{CSS} standard. Some features the actual standard \gls{CSS} 3 offers are not possible. However the \gls{CSS} used in JavaFX supports additional features such as variables for colors. The control part of the previously mentioned \gls{MVC} pattern consists of controller classes specified in the \gls{FXML} files. These controller classes are automatically instantiated when the \gls{FXML} is parsed. With annotations it is possible to link objects to the corresponding \gls{XML} tags to modify them on runtime.

\subsection{Startup Configuration of RoboNav Overview}
\label{subsec:startupconfigurationofrobonavoverview}

When starting \textit{RoboNav Overview} the connection setup screen appears, shown in Figure \ref{fig:connectionsetup}. In this dialogue the user can either connect to an already running instance of \textit{RoboNav Control} by providing the \gls{IP} and the port of the computer running the control software, or starting an own instance on the local system. As explained in section \ref{sec:rncp} an incoming port must be specified which tells the \textit{RoboNav Control} where to send the \textit{\gls{RNCP}} packets and \textit{RoboNav Overview} where to receive them.

\insertpicture{images/connectionsetup.png}{Connection Setup}{(selfmade)}{fig:connectionsetup}{0.4}

After that it is possible to proceed directly to the main window or setup the \textit{RoboNav Configuration} for configuring \textit{RoboNav Control}. The latter choice leads to the configuration setup. In this setup all entries of the \textit{RoboNav Configuration} can be easily customized. These include the connection information for the external camera, the dimensions of the map, the clipping values for the camera image and flags for image modification. The border clipping of the camera image can be adjusted intuitively by dragging the blue borders of the viewport directly on the shown image (Figure \ref{fig:configurationsetup}). When proceeding to the main screen the configuration is automatically sent to the connected \textit{RoboNav Control} instance. It is also possible to save and load configurations to generate \textit{RoboNav Configurations} for \textit{RoboNav Control}-only \textit{RoboNav} systems.

\insertpicture{images/configurationsetup.png}{Configuration Setup}{(selfmade)}{fig:configurationsetup}{0.7}

\subsection{Main Window of RoboNav Overview}

In the main window of \textit{RoboNav Overview} (Figure \ref{fig:mainwindow}), the center is used to display the area map. In the background the overhead camera image is displayed. All connected robots are drawn according to their position. The green circle displays the position where the overhead camera detects the robot. This position is also used for the path calculation. In contrast the red circle shows the position the robot is located according to its internal odometry. The distance of which \textit{RoboNav} improves the position is shown as an orange line with the difference on both axis.
The obstacles are shown in red with the opacity indicating its priority. A higher priority leads to a higher opacity and vice versa.

If the user selects a robot by clicking on it the robot's data is displayed in the right panel. On default this panel shows the robot with the highest priority. In the robot panel the user can see the robot's exact position including its rotation, all sensor values regardless if they are digital or analog and the current battery state. Furthermore the robot's priority can be changed. For convenience the robot can also be navigated manually. If the robot uses an onboard-camera, the image is received by \textit{RoboNav Overview} and displayed at the top of the robot panel.

\insertpicture{images/mainwindow.png}{Main Window}{(selfmade)}{fig:mainwindow}{1}

On the left side of the window, the tool-bar is located. There the tools can be switched. By default the selection tool, which is used to select robots and paths is used. The other tool available is the obstacle tool. With this tool it is possible to create new obstacles manually and send them to \textit{RoboNav Control} so they can be used for avoidance.

The user can create a new path for the selected robot by clicking on the map with the selection tool (Figure \ref{fig:pathselection}) Simultaneously the bottom panel, called path panel, is opened where every point clicked by the user is added to a chronological list. Additionally the user can choose one of the provided algorithms from \textit{RoboNav Control} to travel this path. When the user sends a path-request to \textit{RoboNav Control}, the calculated path is returned to \textit{RoboNav Overview} and automatically displayed on the area map.

\insertpicture{images/createpath.png}{Selecting a path in \textit{RoboNav Overview}}{(selfmade)}{fig:pathselection}{1}

The menu bar provides direct links to the documentation as well as this diploma thesis. Additionally the setup windows mentioned in section \ref{subsec:startupconfigurationofrobonavoverview} can be reopened from the menu.



\newpage %----------------------------------------------------------------------------------------------
\section{Conclusio and Outlook}
\label{sec:conclusio}
\cfoot{Hannah Siegel}
\textit{RoboNav} is designed to enrich the positioning and the navigation of a mobile robot with information from an overhead camera and therefore improving the navigation. The final version of the project showed, that this theoretical idea can be accomplished, but maybe not as easy and well-performing as it was expected. \\ \\
The project should work real-time and show a combination of external tracking in terms of positioning and a map-based external navigation possibility to perform ideal movements. \\
The overall system works fine. The path-planning and the path-driving functionality in particular work just as expected. With well-placed obstacles the paths are driven without collisions. Sometimes the path can not be calculated since no possible way was found, but this is not a bug in the path-planning software, it is only an impossible use-case which is simply ignored by the system. The software has the limitation that before starting the path-driver should be started only when having cleaned the  \texttt{RoboNavPath.txt} file before, otherwise the old path will be driven immediately. Furthermore, the robot must always be started facing in a certain direction, as showed in Figure \ref{fig:mirrorx} , in order for the coordinate systems to match.\\ %TODO recheck
The image processing works very accurate. 
It recognizes even bright coloured objects on white ground, but for good results it's better if the obstacles have a easy distinguish color.
\\
The position comparison on the other hand did not deliver the desired results: even trough the position-information from the overhead-camera can be used from time to time in order to gain accuracy, it can not be used in near real time. The reason is mainly the delay of the image-transmission of the streaming-app and the intensive calculations afterwards. \\ This means, that a functioning position comparison should only be carried out when the robot is resting on the same position for at least two seconds but the average case is still faster than half a second. The fact that a comparison is even available, should theoretically improve any navigation already. The outcome of the position comparison can be seen in Figure \ref{fig:bla}. Since the internal odometry of the \textit{Robotino} is of a high quality, the aberrations were hardly measurable after smaller distances. This made it very hard to compare the improvements with and without \textit{RoboNav} on the small testing area the project had available.
\begin{figure}[htb!]
	\begin{minipage}{.3\textwidth}
		\centering\includegraphics[width=0.8\textwidth]{images/abberation1small}
	\end{minipage}
	\begin{minipage}{.3\textwidth}
		\centering\includegraphics[width=0.9\textwidth]{images/abberation2small}
	\end{minipage}
	\begin{minipage}{.3\textwidth}
		\centering\includegraphics[width=0.9\textwidth]{images/abberation3small}
	\end{minipage}
	\caption[Bla (self-made)]{Aberration at start, after 5 minutes and after 25 minutes}
	\label{fig:bla}
\end{figure}
\FloatBarrier

\textbf{Lessons Learned}\\
The overall idea of the project is fine and the implementation was good enough to ensure an improvement.
For sure, the performance and the speed of processing information could be faster. \textit{Robonav} used Java and therefore ran into some performance and also some parsing problems between the wrapped Java classes and the native C++ library. Maybe using C++ could show some improvements in terms of performance and compatibility issues. Furthermore, even trough the idea of dividing the system into three independent components - \textit{RoboNav Control, RoboNav Overview} and \textit{RoboNav Sight} - is providing a lot of flexibility, it lead to much communication between the components. \\
When working with different systems, such as robots and cameras, the issue of adjusting and scaling their respective coordinates systems should have been considered earlier in the project, since this posed some problems in the final stage of integrating the components, that could not finally be resolved. 
\\\\
%\subsubsection{Ok Conclusio}
%\textbf{Final Outcome changes:} \\
%The overall system works fine. The path-planning and the path-driving functionality in particular work well enough to ensure a stable testing environment but may sometimes fail to execute perfectly. With well-placed obstacles the paths are mostly driven without major collisions. Sometimes the path can not be calculated since no possible way was found, but this is not a bug in the path-planning software, it is only an impossible use-case. The path-driver should be started only when having cleaned the  \texttt{RoboNavPath.txt} file before, otherwise the old path will be driven immediately. Furthermore, the robot must always be started facing XXXX, in order for the coordinate systems to match.\\ 
%\subsubsection{Bad Conclusio}
%\textbf{Final Outcome changes:} \\
%The project should work real-time and show a combination of external tracking in terms of positioning and a map-based external navigation possibility to perform ideal movements. \\
%The path-planning and the path-driving functionality in particular did not work as expected, due to compatibility issues between the coordinate systems of the entities. %With well-placed obstacles the paths are mostly driven without major collisions. Sometimes the path can not be calculated since no possible way was found, but this is not a bug in the path-planning software, it is only an impossible use-case.
%The path-driver should be started only when having cleaned the  \texttt{RoboNavPath.txt} file before, otherwise the old path will be driven immediately. Furthermore, the robot must always be started facing XXXX, in order for the coordinate systems to match.\\ \\
%\textbf{other  changes:} \\
% The fact that a comparison is even available, should theoretically improve any navigation already. Since the internal odometry of the \textit{Robotino} is very high-class, the aberrations of it were hardly measurable after smaller distances. Furthermore, the fact that the navigation component of the system failed to work absolutely fine without eventual bugs, it is not possible to test it over an longer timespan. This made it very hard to compare the improvements with and without \textit{RoboNav}.
\textbf{Outlook}\\
Since the demand for robots is generally growing, their navigation-possibilities will improve as well. Maybe vision based systems are a technology which will get used even more in the future. \\
The \textit{RoboNav} project could be expanded by the following ideas:
\begin{enumerate}
\item \textit{RoboNav} was tested with only one robot, but it was designed for more than that. The usage of more than one robot generates issues such as crossing pathways, priorities or different hardware and software requirements. \\ Whenever using more than one robot, topics like task sharing or intelligent systems may be explored. 
\item  There was only one overhead camera in use. By combining more images, either from a wider range or from different perspectives, more information might have been provided.
\item By programming the \textit{RoboNav Control} in order to gain "\textit{artificial intelligence}", it could take over tasks such as classifing and checking obstacles and using fuzzy sets for handling these obstacles. This would improve the overall navigation due to cutting out eventual errors from the image processing.   
\item \textit{RoboNav} could be integrated into other programming languages, therefore enabling an easier handling of the functions. Blocks for the \textit{RobotinoView} program could help an user to evade the usage of files.
\item The overview camera might be exchanged for a high-resolution industrial camera, therefore delivering faster and better pictures than the android smartphone does. The transmission of the pictures could be implemented in a faster way as well, for example by using two connections - one for the \textit{RoboNav Control} and one for the \textit{RoboNav Overview}. 
\end{enumerate} 


\newpage %----------------------------------------------------------------------------------------------
\cfoot{}
\pagenumbering{roman}
\section{Appendix}
\label{sec:appenix}

\subsection{Glossaries}
\label{subsec:glossaries}
\begingroup
\renewcommand{\section}[2]{}
\printglossary[style=tree]
\endgroup
\newpage

{\small\color{white}.}
\vspace{-2cm}
\subsection{Figures}
\label{subsec:figures}
\begingroup
\renewcommand{\section}[2]{}
\listoffigures
\endgroup

\subsection{Listings}
\label{subsec:listings}
\begingroup
\renewcommand{\section}[2]{}
\lstlistoflistings
\endgroup

\subsection{Sources}
\label{subsec:sources}
\begingroup
\renewcommand{\section}[2]{}
\bibliographystyle{alpha}
\bibliography{sources}
\endgroup

\newpage
\label{subsec:scrumcontract}
\cfoot{Appendix - Scrum contract}
\includepdf[scale=0.75,pages=1,pagecommand=\subsection{Scrum contract}]{documents/scrumVertag}
\includepdf[scale=0.75, pages={2-6}, pagecommand=\texttt{}]{documents/scrumVertag} 

\subsection{UML}
\insertpicture{images/uml-as-svg}{Overall UML}{(selfmade)}{itm:pic1}{1.0}

\subsection{Installation manual}
\label{subsec:isntalman}
\cfoot{Vennesa Belinić}

\subsubsection{Requirements}
Running the \textit{RoboNav} system is only possible with the following software requirements:
\begin{itemize}
	\item Java SDK 8u45 for compiling and running the software. \cite{gitdown}
	\item Git to pull the repository. (It is also possible to download the git-repository as a zip \cite{robonavgit})
	\item Apache Ant to build and run the project. \cite{antdown}
	\item The \texttt{rec::robotino::api2} for controlling the \textit{Robotino v3}. \cite{api2down}
\end{itemize}

\textbf{Installing rec::robotino::api2}\\
Download the install-file for the API from \cite{api2down} and run it. After running it the language has to be choosen. All configurations like installtion location can be left to default. It is important to select the checkbox \textbf{Add application directory to system path} (see Figure \ref{fig:install}).

\begin{figure}[htb!]
	\begin{minipage}{.5\textwidth}
		1.Start the setup
		\centering\includegraphics[width=0.9\textwidth]{images/install1}
	\end{minipage}
	\begin{minipage}{.5\textwidth}
		2.Select the install location
		\centering\includegraphics[width=0.9\textwidth]{images/install2}
	\end{minipage}
	\begin{minipage}{.5\textwidth}
		3.Add application directory to system path
		\centering\includegraphics[width=0.9\textwidth]{images/install3}
	\end{minipage}
	\begin{minipage}{.5\textwidth}
		4.Finish the setup
		\centering\includegraphics[width=0.9\textwidth]{images/install4}
	\end{minipage}
	\caption[Installing rec::robotino::api2 (self-made)]{Installing rec::robotino::api2}
	\label{fig:install}
\end{figure}
\FloatBarrier

\subsubsection{Deploying the RoboNav System}
To deploy the application the git-repository has to be cloned:
\begin{lstlisting}[label=clonerepo, caption=Cloning the git-repository]
$ git clone https://github.com/TGM-HIT/RoboterNavigation
\end{lstlisting}
Before the sources can be compiled the values in the \texttt{build.properties} have to be set correctly (absolute paths are required without \texttt{\textbackslash} ):
\begin{lstlisting}[label=setproperties, caption=Setting the properties]
# the installation directory of the rec::robotino::api2 (use the short path name used by windows, without spaces)
rec.api.robotino3.location = C:/PROGRA~1/RECGMB~1/robotino/api2/bin

# the path to the OpenCv directory
opencv.nativ_libary.location = C:/Users/user/RoboterNavigation/applications/lib/opencv/build/java/x64
\end{lstlisting}
When the repository is cloned successfully and the values of the properties file are set the applications can be compiled:
\begin{lstlisting}[label=compileapps, caption=Compiling the applications]
# run the default target of ant to compile the sources, run the tests, create the javadoc and build the jars
$ ant
\end{lstlisting}

\subsubsection{Running/Using the RoboNav System}
For running and using the application the following ant-target are provided:
\begin{itemize}
	\item \textbf{run-controller}\\runs the \texttt{RoboNavController} by executing the generated jar
	\item \textbf{run-overview}\\starts the \texttt{RoboNav Overview} by running the jar generated for it
	\item \textbf{jar-api}\\generates a non-executable jar with the whole \textit{RoboNav API} packed in 
\end{itemize}
\end{document}


\end{document}
